\definecolor{ImperialBlue}{RGB}{0, 0, 255}    % Imperial Blue
\definecolor{ProcessBlue}{RGB}{0, 145, 212}   % Process Blue for contrast
\definecolor{LightBlue}{RGB}{235, 245, 255}   % Adjusted background blue
\definecolor{CoolGrey}{RGB}{158, 161, 162}    % Neutral Grey

\section{Implementation \& Architecture}


\subsection{High-Level System Design}
MERIT is partitioned into three distinct functional domains: 
\textbf{Extraction}, \textbf{Configuration}, and \textbf{Execution}. 
Figure \ref{fig:high_level_arch} illustrates the orchestration between these layers.

% Figure 4.1: High-Level System Design
\begin{figure}[htbp]
    \centering
    \begin{tikzpicture}[
        node distance=1.2cm and 1.8cm,
        arch_domain/.style={draw=ImperialBlue, fill=LightBlue!30, dashed, rounded corners, thick, inner sep=15pt},
        arch_logic/.style={rectangle, draw=ImperialBlue, fill=white, text width=2.8cm, align=center, minimum height=1.1cm, rounded corners, thick, font=\footnotesize\sffamily, drop shadow},
        arch_user/.style={ellipse, draw=ImperialBlue, fill=ImperialBlue!10, minimum width=1.8cm, minimum height=0.8cm, font=\small\bfseries\sffamily, thick},
        arrow/.style={-{Stealth[scale=1.0]}, thick, draw=ImperialBlue},
        domain_label/.style={font=\footnotesize\bfseries\sffamily, ImperialBlue, anchor=center}
    ]
        % I. Extraction
        \node [arch_logic] (db1) {Supabase\\Unified Store};
        \node [arch_logic, above=0.8cm of db1] (ext) {Multi-source\\Parsing Engine};
        
        % II. Evidence Fusion
        \node [arch_logic, right=2.2cm of db1] (fusion) {Bayesian Fusion\\\& Scoring Engine};
        
        % III. Explainable Interpretation
        \node [arch_logic, right=2.2cm of fusion, yshift=1.6cm] (run) {Candidate Ranker\\\& Scoring Pipeline};
        \node [arch_logic, below=0.5cm of run] (audit) {ShaRP Explanation\\\& Audit Trails};

        % Containers (Using FIT to ensure boxes are inside)
        \begin{scope}[on background layer]
            \node [arch_domain, fit=(ext) (db1)] (area1) {};
            \node [arch_domain, fit=(fusion)] (area2) {};
            \node [arch_domain, fit=(run) (audit)] (area3) {};
        \end{scope}

        % User (Recruiter) - Increased height for absolute clearance
        \node [arch_user, above=5.5cm of area2] (u2) {Recruiter};
        
        % Flows (Segmented for precise label control)
        % To Extraction
        \draw [arrow] (u2.west) -- (u2.west -| area1.center) node[midway, above, font=\footnotesize\sffamily, black, fill=white, inner sep=1pt] {Ingest Documents} -- (ext.north);
        
        % To Configuration
        \draw [arrow] (u2.south) -- (fusion.north) node[pos=0.3, left, font=\footnotesize\sffamily, black] {Set Weights};
        
        % Internal Flows
        \draw [arrow] (ext) -- (db1);
        \draw [arrow] (db1.east) -- (fusion.west);
        \draw [arrow] (fusion.east) -- (run.west);
        \draw [arrow] (run) -- (audit);
        
        % Feedback Loop (Mirrored Symmetry)
        \draw [arrow] (audit.east) -- ++(1.5,0) |- (u2.east) node[pos=0.75, above, font=\footnotesize\sffamily, black, fill=white, inner sep=1pt] {Explainable Insights};

        % Domain Headers (Drawn last with white background to create "gap" in arrows)
        \node [domain_label, fill=white, inner sep=2pt] at (area1.north) {I. DATA EXTRACTION};
        \node [domain_label, fill=white, inner sep=2pt] at (area2.north) {II. MULTI-SOURCE FUSION};
        \node [domain_label, fill=white, inner sep=2pt] at (area3.north) {III. RANKING \& EXPLAINABILITY};

    \end{tikzpicture}
    \caption[MERIT three-stage architecture]{High-Level System Design: The three stages of MERIT: Ingestion, Scoring, Explainability}
    \label{fig:high_level_arch}
\end{figure}


\subsubsection{Architectural Rationale}
MERIT employs a decoupled microservice architecture to separate the 
concerns of data ingestion, job configuration, and the presentation of configuration results.
Flask (Python) was selected as a backend, providing a robust foundation for data 
ingestion and fusion. Next.js was selected for the frontend to enable a high-performance, 
interactive interface for "Glass-Box" visualisations.

A relational PostgreSQL store was selected to enforce data integrity, ensuring that every 
ranked result is linked via foreign keys to a specific, immutable recruiter configuration, 
allowing for a complete audit trail from "Source-to-Score." We opted for Supabase to 
streamline development and to avoid the overhead of setting up and maintaining our own database.
The relational nature of databases facilitates the merging of disparate data sources
into a single, queryable identity hub. This allows the scoring engine to resolve conflicting 
evidence while maintaining a historical record of all data sources.

\subsection{The Extraction Layer}
The Extraction Layer (Figure \ref{fig:extraction_flow}) is responsible for both the ingestion and 
normalisation of heterogeneous data. Rather than a monolithic service, normalisation is 
implemented as a logical pipeline that performs canonical skill mapping and enforces a unified schema. 
This design decision was made to allow for easy extension of the system to support new 
data sources, allowing MERIT to stay updated with the latest tools and platforms.

% Figure 4.2: The Extraction Layer
\begin{figure}[htbp]
    \centering
    \begin{tikzpicture}[
        node distance=1.2cm and 1.5cm,
        arch_source/.style={rectangle, draw=CoolGrey, fill=LightBlue!20, text width=2.4cm, align=center, minimum height=0.8cm, font=\footnotesize\sffamily, drop shadow},
        arch_logic/.style={rectangle, draw=ImperialBlue, fill=white, text width=2.8cm, align=center, minimum height=1cm, rounded corners, thick, font=\footnotesize\sffamily, drop shadow},
        arch_ir/.style={cylinder, draw=ImperialBlue, fill=LightBlue!40, text width=1.8cm, align=center, minimum height=1.4cm, shape border rotate=90, font=\footnotesize\sffamily, drop shadow},
        field/.style={font=\footnotesize\sffamily, black},
        arrow/.style={-{Stealth[scale=1.0]}, thick, draw=ImperialBlue}
    ]
        % Column 1: Sources (CV/Github/LinkedIn)
        \node [arch_source] (cv) {\textbf{CV (PDF/DOCX)}\\Unstructured Text};
        \node [arch_source, below=1.8cm of cv] (api) {\textbf{GitHub/LinkedIn}\\Structured API};

        % Column 2: Specific Engines (The parts on the right of the sources)
        \node [arch_logic, right=of cv] (nlp) {Heuristic NLP Engine\\{\footnotesize pdfplumber + regex}};
        \node [arch_logic, right=of api] (agg) {API Aggregator\\{\footnotesize REST / GraphQL}};

        % Column 3: Normaliser (The combination block)
        \node [arch_logic, right=1cm of nlp, yshift=-1.4cm, text width=2.5cm, fill=LightBlue!10] (norm) {Normalisation\\Engine\\{\footnotesize (Logic Layer)}};
        
        % Column 4: The Database
        \node [arch_ir, right=of norm] (output) {Unified\\Candidate Store\\{\footnotesize Postgres / Supabase}};

        % Labels
        \node [field, above=0.1cm of nlp] {Extracts Work, Skills, Education};
        \node [field, below=0.1cm of agg] {Extracts Commits, Stars, Badges, Languages (Python, Java, etc.)};

        % Connections
        \draw [arrow] (cv) -- (nlp);
        \draw [arrow] (api) -- (agg);
        
        % Orthogonal arrows to normalizer
        \draw [arrow] (nlp.east) -- ++(0.5,0) |- (norm.west);
        \draw [arrow] (agg.east) -- ++(0.5,0) |- (norm.west);
        \draw [arrow] (norm) -- (output);

    \end{tikzpicture}
    \caption{Stage 1: The Intelligent Extraction Pipeline, decomposing unstructured 
    documents and structured API signals into a unified, queryable schema.}
    \label{fig:extraction_flow}
\end{figure}


MERIT is designed to be able to ingest both structured and unstructured data sources.
For unstructured data sources, such as CVs and JDs, we employ \textit{pdfplumber} as well
as a custom heuristic NLP engine to extract key entities such as work experience, education 
history, and technical skills.

For structured data sources, such as GitHub and LinkedIn, we leverage APIs to extract
meta-data such as commit frequencies, repository stars, and programming language 
distributions. By unifying these disparate signals into a single schema, MERIT is able to provide
comprehensive views of a candidate's technical profile, with verifiability and transparency at its core


\subsection{The Configuration and Execution Layer}
% Figure 4.3: The relational Schema
\begin{figure}[htbp]
    \centering
    \begin{tikzpicture}[
        node distance=1.5cm and 2.5cm,
        arch_table/.style={rectangle, draw=CoolGrey, fill=white, text width=3.2cm, align=left, font=\footnotesize\sffamily, rounded corners=1pt, drop shadow, minimum height=1.6cm},
        arch_source/.style={ellipse, draw=ImperialBlue, fill=LightBlue!40, font=\footnotesize\bfseries\sffamily, minimum width=2.0cm, align=center},
        arch_group/.style={draw=ImperialBlue, fill=LightBlue!10, dashed, rounded corners, inner sep=10pt},
        pk/.style={font=\footnotesize\bfseries\sffamily, ImperialBlue},
        fk/.style={font=\footnotesize\itshape\color{ProcessBlue}},
        arrow/.style={-{Stealth[scale=1.0]}, thick, draw=ImperialBlue}
    ]
        % Candidate Schema Part: Sources in a column on the left
        \node [arch_table] (t_gh) {
            \textbf{github\_profiles}\\
            {\color{ImperialBlue}\bfseries id}\\
            username, total\_commits\\
            total\_stars, languages
        };

        \node [arch_table, below=0.6cm of t_gh] (t_li) {
            \textbf{linkedin\_profiles}\\
            {\color{ImperialBlue}\bfseries id}\\
            full\_name, headline\\
            location, about
        };

        \node [arch_table, below=0.6cm of t_li] (t_edu) {
            \textbf{candidate\_education}\\
            {\color{ImperialBlue}\bfseries id}\\
            {\itshape\color{ProcessBlue} candidate\_id}\\
            school\_name, degree, grade
        };

        % Central Identity Table (Closer to sources)
        \node [arch_table, right=1.6cm of t_li] (t_cand) {
            \textbf{candidate\_data}\\
            {\color{ImperialBlue}\bfseries id}\\
            name, email, skills\\
            {\itshape\color{ProcessBlue} github\_profile\_id}\\
            {\itshape\color{ProcessBlue} linkedin\_profile\_id}
        };

        % Configuration Schema Part: Anchored further right to avoid collision
        \node [arch_table, right=6cm of t_gh] (t_jd) {
            \textbf{job\_requirements}\\
            {\color{ImperialBlue}\bfseries id}\\
            title, description\\
            metrics (JSONB)
        };

        \node [arch_table, below=0.6cm of t_jd] (t_conf) {
            \textbf{matching\_configs}\\
            {\color{ImperialBlue}\bfseries id}\\
            {\itshape\color{ProcessBlue} job\_id}, {\itshape\color{ProcessBlue} batch\_id}\\
            weights, active\_metrics
        };

        % Results
        \node [arch_table, below=0.6cm of t_conf] (t_res) {
            \textbf{past\_results}\\
            {\color{ImperialBlue}\bfseries id}\\
            {\itshape\color{ProcessBlue} config\_id}\\
            results\_payload (JSONB)
        };

        % Groupings
        \begin{scope}[on background layer]
            \node [arch_group, fit=(t_gh) (t_li) (t_edu) (t_cand), label={[font=\footnotesize\bfseries\sffamily, ImperialBlue]above:CANDIDATE PERSISTENCE}] (db_cand) {};
            \node [arch_group, fit=(t_jd) (t_conf) (t_res), label={[font=\footnotesize\bfseries\sffamily, ImperialBlue]above:CONFIG \& RESULTS}] (db_conf) {};
        \end{scope}

        % Flows
        \draw [arrow] (t_jd) -- (t_conf);
        \draw [arrow] (t_conf) -- (t_res);
        
        % Internal Relationships
        \draw [arrow] (t_gh.east) -- ++(0.4,0) |- (t_cand.west);
        \draw [arrow] (t_li.east) -- (t_cand.west);
        \draw [arrow] (t_edu.east) -- ++(0.4,0) |- (t_cand.west);

        % Flow to Matching (Clear crossing between groups)
        \draw [arrow] (t_cand.east) -- (t_conf.west);

    \end{tikzpicture}
    \caption{Relational Schema: Integration of candidate signals, recruiter criteria, 
    and auditable matching results.}
    \label{fig:schema_overview}
\end{figure}

The core of MERIT's "Glass-Box" transparency lies in its relational 
data model (Figure \ref{fig:schema_overview}). This model bridges the 
\textbf{Configuration} and \textbf{Execution} layers by ensuring every extracted 
signal is persisted in a way that remains queryable by the scoring engine.


\subsection{The Presentation Layer}
The Presentation Layer is the user interface of MERIT, providing an interactive 
and seamless experience for recruiters to evaluate candidates. It is built 
using Next.js, which enables high-performance, dynamic web applications. 
We fetch data from the previously mentioned Flask API, and reflect it
in a dynamic, interactive dashboard. We provide explainability through audit trails
that trace the results back to the specific recruiter configurations. This makes 
the decision-making process auditable and verifiable.
